% !TeX root = ../main.tex

% 中文摘要
\begin{abstract}
  本文关注概率模型在中学恒等式证明中的应用，通过文献研究法、对比研究法和案例分析法，研究其证明方式。文章阐释了概率模型的相关理论基础，包括概率模型基本概念、概率模型基本类型和适用原理。详细介绍利用构造等可能事件概率模型、独立重复试验概率模型等方法，并给出相应实例和分析。还提出了该方法在中学教学中的应用建议，探讨这种方法对培养学生思维能力的作用，为中学数学教学与学习提供新的思路，帮助学生提升数学素养与综合能力。
\end{abstract}



% 英文摘要
\begin{abstract*}
  This article focuses on the application of probability models in the proof of identities in middle schools. Through the methods of literature research, comparative research and case analysis, the proof methods are studied. The article explains the relevant theoretical basis of probability models, including the basic concepts of probability models, the basic types of probability models and the applicable principles. The methods such as constructing the probability model of equal possible events and the probability model of independent repeated trials are introduced in detail, and corresponding examples and analyses are given. It also put forward application suggestions of this method in middle school teaching, explored the role of this method in cultivating students' thinking ability, provided new ideas for middle school mathematics teaching and learning, and helped students improve their mathematical literacy and comprehensive ability.

  % This article is dedicated to the application of probability models in proving identities in high school mathematics. Through literature review, comparative studies, and case analysis, it examines the methods used for proof. The article explains the theoretical foundations of probability models, including basic concepts, fundamental types, and applicable principles. It provides detailed introductions to methods such as constructing equal possible events and independent repeated trials probability models, along with corresponding examples and analyses. It gives advice on applying these methods in high school teaching, exploring their role in fostering students' thinking skills. This provides new perspectives for high school mathematics instruction and learning, helping students enhance their mathematical literacy and comprehensive abilities.
\end{abstract*}